\section{Introduction}


意識をめぐる議論はしばしば、「どこで生成されるのか」「どの回路が担うのか」「どの表象形式に対応するのか」という問いから始まる。本稿は、この問いの置き方自体を変更する。意識は本稿ではどこかで作られる対象として扱われず、異なる時間スケールの処理が、一定の遅延と結合を保ったまま相互参照可能になるときに開かれる状態としてモデル化される。

この見方では、意識は自己、思考、判断そのものではない。意識は、高速な予測過程が低速な記憶・報告・沈殿の層へ参照可能になる条件である。したがって、本稿の目的は意識の最終定義を与えることではなく、意識を論じるための構造的位置を明確にすることである。

既存の補遺的議論では、クオリア、Libet 実験、分離脳、測定、局在化問題が個別に扱われていた。本稿ではそれらを一つの独立した理論草稿として再編成する。

\subsection{本稿の立場：三つの宣言}

本節では、本稿が既存の意識理論群に対してどのような立場を取るかを、三点に絞って先に示す。詳細な比較は §18 で行う。ここでは読者が論文全体を通じて保持すべき視点の座標を固定することを目的とする。

\textbf{第一：生成モデルではなく開口モデル}

本稿は、意識を「どこかで生成される出力」として扱わない。Global Workspace 型・IIT 型・Higher-Order 型を含む多くのモデルは、意識をある処理の産物、あるいは統合量の到達点として記述する。本稿はその前提を出発点として採用しない。

より正確には、本稿では意識を生成された対象としてモデル化しない。意識は、高速な予測過程と低速な沈殿的記憶過程のあいだに、特定の速度差・結合粘性・内部遅延・意味勾配が成立したときに開かれる状態（ログ参照状態）として扱われる。この違いは定義上の差異ではなく、意識について何を問うべきかという問いの構造そのものに関わる。

\textbf{第二：単一量への還元ではなく時間的レジームとしての記述}

本稿は、意識を単一の物理量・情報量・統合度に帰着させない。意識は、速度差 $\Delta v$、結合粘性 $\mu$、内部遅延 $\tau$、有効トルク様量 $\mathcal{T}_{\Phi}$、意味勾配 $\nabla\Phi$ のいずれか一つに還元されるものではない。意識は、それらの変数が一定の関係に入ったときに成立する時間的非閉包レジームである。

この立場は消去主義ではない。意識を物理量や神経活動の記述に「溶かす」ことを目指していない。また二元論でもない。意識という特別な実体を追加することもしない。意識を、複数の処理層の動的関係として記述する。

\textbf{第三：神経科学の否定ではなく条件の再記述}

本稿は、意識が単一部位に局在しないと主張する。しかしこれは神経科学の否定ではない。ある脳部位・回路・振動・活動パターンが意識状態と相関することは十分にありうる。本稿の立場では、それらはログ参照ループに参加する局所的条件として解釈される。

神経相関（neural correlates of consciousness）は意識そのものではない。それは、意識が開かれるための局所基盤の一部である。「意識の座を探す」問いから「意識が開かれる条件を調べる」問いへの移行が、本稿の神経科学に対するスタンスである。

これら三点を保持した上で、以下の論述を読んでほしい。本稿が主張しているのは、意識を否定することでも神秘化することでもなく、意識について問える構造的位置を固定することである。

なお、本稿で扱う高速層・低速層の二層構造は、IFGT（Information Field Geometry Theory）、Bio-IFGT、および Conceptual Topography における認知的観測窓の射影と接続している。IFGT は情報密度・情報流・情報ポテンシャルの一般語彙を与え、Bio-IFGT はそのような層構造が生物システムに現れうる理由を示し、CT は概念地形が認知的観測窓の中でどのように安定化するかを記述する。なぜこの二層が生物システムに現れるのかという問いへの根拠は §4 で示す。本稿を独立した論考として読む場合も、その接続を知ることで式の意味が深まる。
